\subsection{Conditional correspondence of the two strategies}
\label{sec:mode-correspondence}
Sound summaries alone do not imply equal test results:
$\{1\}\subseteq\{0,1\}$ is a sound abstraction, but the additional result
0 refutes a caller's requirement that the return be 1.
Likewise, an inductive invariant may contain states unreachable at a
particular iteration. Conversely, a single summary sample may omit an
implementation behavior. We therefore state a sufficient correspondence
condition explicitly rather than assume agreement from a mode switch.

Let $i\leadsto_I i'$ and $m\leadsto_M m'$ be steps between corresponding
boundaries of body execution and summary execution.
A boundary may package several C steps only if their safety conditions
and selected observations are retained; an endpoint-only summary cannot
silently hide intermediate violations.
Let $\mathcal R(i,m)$ relate their states, including immutable bindings,
live values, visible heap resources and frames.
The interface requires:
\begin{enumerate}
\item the initial domains cover each other through $\mathcal R$;
\item related states have equal observation views and agree on safety,
selected checks (including return checks), and terminality;
\item every body successor has a related summary successor, and every
summary successor has a related body successor.
\end{enumerate}
The last condition is a local two-way simulation:
\[
\begin{split}
\mathcal R(i,m)\land i\leadsto_I i'
 &\Rightarrow \exists m'.\ m\leadsto_M m'\land\mathcal R(i',m'),\\
\mathcal R(i,m)\land m\leadsto_M m'
 &\Rightarrow \exists i'.\ i\leadsto_I i'\land\mathcal R(i',m').
\end{split}
\]
An observation view records selected check identities and order, visible
values and resources. The same schedule and interpretation are used on both
sides. Any additional modular preservation obligations must also be
justified; matching final return values alone does not discharge them.

\begin{theorem}[Conditional strategy correspondence]
\label{thm:mode-correspondence}
Under this interface, related states have matching finite boundary traces
with the same observation views. They agree on existence of a reachable
violation, on correctness of complete finite traces, and on semantic
well-founded termination. Nonvacuous total acceptance agrees between the
related initial domains when acceptance requires safety and all selected
checks, progress at nonterminal states, and termination without outgoing
steps at terminal states.
\end{theorem}
\begin{proof}
Induction on a finite trace applies the appropriate successor clause at
each step and preserves $\mathcal R$; the observation and check clauses
transfer the recorded views and their truth. This transports both a
violating endpoint and every complete trace in either direction.
Induction on the total-acceptance derivation transports safety, checks,
terminality, and progress; the reverse successor clause ensures that
\emph{every} successor on the other side is covered by the induction
hypothesis. Initial-domain coverage transfers nonemptiness and the
universal acceptance condition. The same successor matching transports
well-founded termination.
\end{proof}

The standard-library Rocq module \texttt{ModeCorrespondence.v} proves
these relational claims and the loose-summary counterexample.
The interface is an obligation for a concrete summary construction,
not a claim that arbitrary contracts, invariants, or the current sampler
satisfy it. In particular, choosing one legal model is insufficient.
Termination here concerns the supplied boundary transition systems.
Lifting it to C also requires finite boundary segments that do not hide
internal divergence or stuck behavior. The theorem compares semantic outcomes. Identical finite-resource tool
reports additionally require complete execution coverage and successful
decision of corresponding obligations; different solver support, proof
search, or time limits may leave one report UNKNOWN.
